\documentclass{amsart}
% For dealing with references we use the comment environment
\usepackage{verbatim}
\newenvironment{reference}{\comment}{\endcomment}
%\newenvironment{reference}{}{}
\newenvironment{slogan}{\comment}{\endcomment}
\newenvironment{history}{\comment}{\endcomment}

% For commutative diagrams we use Xy-pic
\usepackage[all]{xy}

% We use 2cell for 2-commutative diagrams.
\xyoption{2cell}
\UseAllTwocells

% We use multicol for the list of chapters between chapters
\usepackage{multicol}

% This is generally recommended for better output
\usepackage{lmodern}
\usepackage[T1]{fontenc}

% For cross-file-references

% Package for hypertext links:
\usepackage{hyperref}

% For any local file, say "hello.tex" you want to link to please

% Theorem environments.
%
\theoremstyle{plain}
\newtheorem{theorem}[subsection]{Theorem}
\newtheorem{proposition}[subsection]{Proposition}
\newtheorem{lemma}[subsection]{Lemma}

\theoremstyle{definition}
\newtheorem{definition}[subsection]{Definition}
\newtheorem{example}[subsection]{Example}
\newtheorem{exercise}[subsection]{Exercise}
\newtheorem{situation}[subsection]{Situation}

\theoremstyle{remark}
\newtheorem{remark}[subsection]{Remark}
\newtheorem{remarks}[subsection]{Remarks}

\numberwithin{equation}{subsection}

% Macros
%
\def\lim{\mathop{\mathrm{lim}}\nolimits}
\def\colim{\mathop{\mathrm{colim}}\nolimits}
\def\Spec{\mathop{\mathrm{Spec}}}
\def\Hom{\mathop{\mathrm{Hom}}\nolimits}
\def\Ext{\mathop{\mathrm{Ext}}\nolimits}
\def\SheafHom{\mathop{\mathcal{H}\!\mathit{om}}\nolimits}
\def\SheafExt{\mathop{\mathcal{E}\!\mathit{xt}}\nolimits}
\def\Sch{\mathit{Sch}}
\def\Mor{\mathop{\mathrm{Mor}}\nolimits}
\def\Ob{\mathop{\mathrm{Ob}}\nolimits}
\def\Sh{\mathop{\mathit{Sh}}\nolimits}
\def\NL{\mathop{N\!L}\nolimits}
\def\CH{\mathop{\mathrm{CH}}\nolimits}
\def\proetale{{pro\text{-}\acute{e}tale}}
\def\etale{{\acute{e}tale}}
\def\QCoh{\mathit{QCoh}}
\def\Ker{\mathop{\mathrm{Ker}}}
\def\Im{\mathop{\mathrm{Im}}}
\def\Coker{\mathop{\mathrm{Coker}}}
\def\Coim{\mathop{\mathrm{Coim}}}

% Boxtimes
%
\DeclareMathSymbol{\boxtimes}{\mathbin}{AMSa}{"02}

%
% Macros for moduli stacks/spaces
%
\def\QCohstack{\mathcal{QC}\!\mathit{oh}}
\def\Cohstack{\mathcal{C}\!\mathit{oh}}
\def\Spacesstack{\mathcal{S}\!\mathit{paces}}
\def\Quotfunctor{\mathrm{Quot}}
\def\Hilbfunctor{\mathrm{Hilb}}
\def\Curvesstack{\mathcal{C}\!\mathit{urves}}
\def\Polarizedstack{\mathcal{P}\!\mathit{olarized}}
\def\Complexesstack{\mathcal{C}\!\mathit{omplexes}}
% \Pic is the operator that assigns to X its picard group, usage \Pic(X)
% \Picardstack_{X/B} denotes the Picard stack of X over B
% \Picardfunctor_{X/B} denotes the Picard functor of X over B
\def\Pic{\mathop{\mathrm{Pic}}\nolimits}
\def\Picardstack{\mathcal{P}\!\mathit{ic}}
\def\Picardfunctor{\mathrm{Pic}}
\def\Deformationcategory{\mathcal{D}\!\mathit{ef}}

\setlength{\emergencystretch}{3em}
\begin{document}
\title[Stacks Project, Tag 06SM]{398 --- Existence of a smooth complete local base with prescribed cotangent sequence}
\maketitle
\noindent Copyright (C) 2005--2025 Johan de Jong. Stacks Project authors. Source: \href{https://stacks.math.columbia.edu/tag/06SM}{Tag 06SM}.

\noindent Editorial difficulty note (not author text). Difficulty H2: For a possibly inseparable residue extension, the proof successively slices a smooth algebra to eliminate nonzero differential classes. Completing the resulting local algebra preserves its cotangent data and formal smoothness. This constructs a base whose cotangent sequence has the asserted two isomorphisms, beyond the simpler separable finite-etale construction..

\noindent Editorial provenance note (not author text). History: excerpt from revision \texttt{a0444\allowbreak 6e57e\allowbreak c1fbc\allowbreak 252a8\allowbreak 71afc\allowbreak ec775\allowbreak 2fb28\allowbreak 07b14}, formal-defos.tex, lines 2380--2464. Reference presentation was adapted; original-author context and core support blocks were retained or added. The mathematical wording is unchanged. Licensed under GNU FDL 1.2 or later; no invariant sections or cover texts. See ../licenses/COPYING. Context blocks reproduce author definitions and supporting results; they are not additional counted problems.

\begin{definition}
\label{definition-CLambda}
Let $\Lambda$ be a Noetherian ring and let $\Lambda \to k$ be a finite
ring map where $k$ is a field. We define {\it $\mathcal{C}_\Lambda$} to be
the category with
\begin{enumerate}
\item objects are pairs $(A, \varphi)$ where $A$ is an Artinian local
$\Lambda$-algebra and where $\varphi : A/\mathfrak m_A \to k$ is a
$\Lambda$-algebra isomorphism, and
\item morphisms $f : (B, \psi) \to (A, \varphi)$ are local $\Lambda$-algebra
homomorphisms such that $\varphi \circ (f \bmod \mathfrak m) = \psi$.
\end{enumerate}
We say we are in the {\it classical case} if $\Lambda$ is a Noetherian
complete local ring and $k$ is its residue field.
\end{definition}

\section{Notation and Conventions}
\label{section-notations-conventions}

\noindent
A ring is commutative with $1$. The  maximal ideal of a local ring $A$
is denoted by $\mathfrak{m}_A$. The set of positive integers is denoted
by $\mathbf{N} = \{1, 2, 3, \ldots\}$. If $U$ is an object of a
category $\mathcal{C}$, we denote by $\underline{U}$
the functor
$\Mor_\mathcal{C}(U, -): \mathcal{C} \to \textit{Sets}$, see
Remarks \href{https://stacks.math.columbia.edu/tag/06GK}{remarks cofibered groupoids (Tag 06GK)} (\href{https://stacks.math.columbia.edu/tag/06GQ}{item definition yoneda (Tag 06GQ)}).
Warning: this may conflict with the notation in other chapters where we
sometimes use $\underline{U}$ to denote $h_U(-) = \Mor_\mathcal{C}(-, U)$.

\medskip\noindent
Throughout this chapter $\Lambda$ is a Noetherian ring and
$\Lambda \to k$ is a finite ring map from $\Lambda$ to a field.
The kernel of this map is denoted $\mathfrak m_\Lambda$ and the
image $k' \subset k$. It turns out that $\mathfrak m_\Lambda$ is
a maximal ideal, $k' = \Lambda/\mathfrak m_\Lambda$ is a field, and
the extension $k/k'$ is finite. See discussion surrounding
(\href{https://stacks.math.columbia.edu/tag/06S2}{equation k prime (Tag 06S2)}).




\begin{definition}
\label{definition-completion-CLambda}
Let $\Lambda$ be a Noetherian ring and let $\Lambda \to k$ be a finite
ring map where $k$ is a field. We define {\it $\widehat{\mathcal{C}}_\Lambda$}
to be the category with
\begin{enumerate}
\item objects are pairs $(R, \varphi)$ where $R$ is a Noetherian complete
local $\Lambda$-algebra and where $\varphi : R/\mathfrak m_R \to k$ is a
$\Lambda$-algebra isomorphism, and
\item morphisms $f : (S, \psi) \to (R, \varphi)$ are local $\Lambda$-algebra
homomorphisms such that $\varphi \circ (f \bmod \mathfrak m) = \psi$.
\end{enumerate}
\end{definition}

\begin{definition}
\label{definition-cofibered-groupoid-projection-smooth}
Let $p : \mathcal{F} \to \mathcal{C}_\Lambda$ be a category cofibered in
groupoids. We say $\mathcal{F}$ is {\it smooth} or {\it unobstructed}
if its structure morphism $p$ is smooth
in the sense of Definition \ref{definition-smooth-morphism}.
\end{definition}

\noindent
To analyze essential surjections in $\mathcal{C}_\Lambda$ a bit more
we introduce some notation. Suppose that $A$ is an object
of $\mathcal{C}_\Lambda$ or more generally any $\Lambda$-algebra
equipped with a $\Lambda$-algebra surjection $A \to k$.
There is a canonical exact sequence
\begin{equation}
\label{equation-sequence}
\mathfrak m_A/\mathfrak m_A^2 \xrightarrow{\text{d}_A}
\Omega_{A/\Lambda} \otimes_A k \to
\Omega_{k/\Lambda} \to 0
\end{equation}
see
Algebra, Lemma \href{https://stacks.math.columbia.edu/tag/00RU}{lemma differential seq (Tag 00RU)}.
Note that $\Omega_{k/\Lambda} = \Omega_{k/k'}$ with $k'$ as
in (\href{https://stacks.math.columbia.edu/tag/06S2}{equation k prime (Tag 06S2)}). Let $H_1(L_{k/\Lambda})$
be the first homology module of the naive cotangent complex of $k$
over $\Lambda$, see
Algebra, Definition \ref{algebra-definition-naive-cotangent-complex}.
Then we can extend (\ref{equation-sequence})
to the exact sequence
\begin{equation}
\label{equation-sequence-extended}
H_1(L_{k/\Lambda}) \to
\mathfrak m_A/\mathfrak m_A^2 \xrightarrow{\text{d}_A}
\Omega_{A/\Lambda} \otimes_A k \to
\Omega_{k/\Lambda} \to 0,
\end{equation}
see
Algebra, Lemma \href{https://stacks.math.columbia.edu/tag/00S2}{lemma exact sequence NL (Tag 00S2)}.
If $B \to A$ is a ring map in $\mathcal{C}_\Lambda$
or more generally a map of $\Lambda$-algebras equipped
with $\Lambda$-algebra surjections onto $k$, then we obtain a
commutative diagram
\begin{equation}
\label{equation-sequence-functorial}
\vcenter{
\xymatrix{
H_1(L_{k/\Lambda}) \ar[r] \ar@{=}[d] &
\mathfrak m_B/\mathfrak m_B^2 \ar[r]_{\text{d}_B} \ar[d] &
\Omega_{B/\Lambda} \otimes_B k \ar[r] \ar[d] &
\Omega_{k/\Lambda} \ar[r] \ar@{=}[d] & 0 \\
H_1(L_{k/\Lambda}) \ar[r] &
\mathfrak m_A/\mathfrak m_A^2 \ar[r]^{\text{d}_A} &
\Omega_{A/\Lambda} \otimes_A k \ar[r] &
\Omega_{k/\Lambda} \ar[r] & 0
}
}
\end{equation}
with exact rows.

\begin{definition}
\label{definition-smooth-morphism}
Let $\varphi : \mathcal{F} \to \mathcal{G}$ be a morphism of categories
cofibered in groupoids over $\mathcal{C}_\Lambda$.  We say  $\varphi$ is
{\it smooth} if it satisfies the following condition: Let $B \to A$ be
a surjective ring map in $\mathcal{C}_\Lambda$.  Let $y \in
\Ob(\mathcal{G}(B)), x \in \Ob(\mathcal{F}(A))$, and $y
\to \varphi(x)$ be a morphism lying over $B \to A$.  Then there
exists $x' \in \Ob(\mathcal{F}(B))$, a morphism $x' \to x$
lying over $B \to A$, and a morphism $\varphi(x') \to y$ lying
over $\text{id}: B \to B$, such that the diagram
$$
\xymatrix{
\varphi(x') \ar[r] \ar[dr] & y \ar[d] \\
& \varphi(x)
}
$$
commutes.
\end{definition}

\begin{definition}
\label{algebra-definition-naive-cotangent-complex}
Let $R \to S$ be a ring map. The {\it naive cotangent complex}
$\NL_{S/R}$ is the chain complex (\href{https://stacks.math.columbia.edu/tag/07BM}{equation naive cotangent complex (Tag 07BM)})
$$
\NL_{S/R} = \left(I/I^2 \longrightarrow \Omega_{R[S]/R} \otimes_{R[S]} S\right)
$$
with $I/I^2$ placed in (homological) degree $1$ and
$\Omega_{R[S]/R} \otimes_{R[S]} S$ placed in degree $0$. We will denote
$H_1(L_{S/R}) = H_1(\NL_{S/R})$\footnote{This module is sometimes
denoted $\Gamma_{S/R}$ in the literature.} the homology in degree $1$.
\end{definition}

\noindent
Explicitly, $p: \mathcal{F} \to \mathcal{C}$ is cofibered in groupoids if
the following two conditions hold:
\begin{enumerate}
\item For every morphism $f: U \to V$ in $\mathcal{C}$ and every object
$x$ lying over $U$, there is a morphism $x \to y$ of $\mathcal{F}$ lying
over $f$.
\item For every pair of morphisms $a: x \to y$ and $b: x \to z$
of $\mathcal{F}$ and any morphism $f: p(y) \to p(z)$ such that $p(b) = f
\circ p(a)$, there exists a unique morphism $c: y \to z$ of $\mathcal
F$ lying over $f$ such that $b = c \circ a$.
\end{enumerate}


\begin{lemma}
\label{lemma-exists-smooth}
There exists an $R \in \Ob(\widehat{\mathcal{C}}_\Lambda)$
such that the equivalent conditions of Lemma \href{https://stacks.math.columbia.edu/tag/0DYL}{lemma smooth (Tag 0DYL)}
hold and moreover $H_1(L_{k/\Lambda}) = \mathfrak m_R/\mathfrak m_R^2$
and $\Omega_{R/\Lambda} \otimes_R k = \Omega_{k/\Lambda}$.
\end{lemma}

\begin{proof}
In the classical case we choose $R = \Lambda$. More generally, if
the residue field extension $k/k'$ is separable, then there exists
a unique finite \'etale extension $\Lambda^\wedge \to R$
(Algebra, Lemmas \href{https://stacks.math.columbia.edu/tag/04GM}{lemma complete henselian (Tag 04GM)} and
\href{https://stacks.math.columbia.edu/tag/04GK}{lemma henselian cat finite etale (Tag 04GK)})
of the completion $\Lambda^\wedge$ of $\Lambda$
inducing the extension $k/k'$ on residue fields.

\medskip\noindent
In the general case we proceed as follows. Choose a smooth
$\Lambda$-algebra $P$ and a $\Lambda$-algebra surjection $P \to k$.
(For example, let $P$ be a polynomial algebra.)
Denote $\mathfrak m_P$ the kernel of $P \to k$. The Jacobi-Zariski sequence,
see (\ref{equation-sequence-extended}) and 
Algebra, Lemma \href{https://stacks.math.columbia.edu/tag/00S2}{lemma exact sequence NL (Tag 00S2)}, is an
exact sequence
$$
0 \to H_1(\NL_{k/\Lambda}) \to
\mathfrak m_P/\mathfrak m_P^2 \to
\Omega_{P/\Lambda} \otimes_P k \to
\Omega_{k/\Lambda} \to 0
$$
We have the $0$ on the left because $P/k$ is smooth, hence
$\NL_{P/\Lambda}$ is quasi-isomorphic to a finite projective module
placed in degree $0$, hence $H_1(\NL_{P/\Lambda} \otimes_P k) = 0$.
Suppose $f \in \mathfrak m_P$ maps to a nonzero element
of $\Omega_{P/\Lambda} \otimes_P k$. Setting $P' = P/(f)$ we have a
$\Lambda$-algebra surjection $P' \to k$.
Observe that $P'$ is smooth at $\mathfrak m_{P'}$:
this follows from More on Morphisms, Lemma
\href{https://stacks.math.columbia.edu/tag/057C}{lemma slice smooth given element (Tag 057C)}.
Thus after replacing $P$ by a principal localization
of $P'$, we see that $\dim(\mathfrak m_P/\mathfrak m_P^2)$
decreases. Repeating finitely many times, we may assume
the map $\mathfrak m_P/\mathfrak m_P^2 \to
\Omega_{P/\Lambda} \otimes_P k$ is zero so that
the exact sequence breaks into isomorphisms
$H_1(L_{k/\Lambda}) = \mathfrak m_P/\mathfrak m_P^2$ and
$\Omega_{P/\Lambda} \otimes_P k = \Omega_{k/\Lambda}$.

\medskip\noindent
Let $R$ be the $\mathfrak m_P$-adic completion of $P$.
Then $R$ is an object of $\widehat{\mathcal{C}}_\Lambda$.
Namely, it is a complete local Noetherian ring (see
Algebra, Lemma \href{https://stacks.math.columbia.edu/tag/0316}{lemma completion Noetherian Noetherian (Tag 0316)})
and its residue field is identified with $k$.
We claim that $R$ works.

\medskip\noindent
First observe that the map $P \to R$ induces isomorphisms
$\mathfrak m_P/\mathfrak m_P^2 = \mathfrak m_R/\mathfrak m_R^2$
and $\Omega_{P/\Lambda} \otimes_P k = \Omega_{R/\Lambda} \otimes_R k$.
This is true because both $\mathfrak m_P/\mathfrak m_P^2$
and $\Omega_{P/\Lambda} \otimes_P k$ only depend on the
$\Lambda$-algebra $P/\mathfrak m_P^2$, see
Algebra, Lemma \href{https://stacks.math.columbia.edu/tag/02HQ}{lemma differential mod power ideal (Tag 02HQ)},
the same holds for $R$ and we have $P/\mathfrak m_P^2 = R/\mathfrak m_R^2$.
Using the functoriality of the Jacobi-Zariski sequence
(\ref{equation-sequence-functorial})
we deduce that $H_1(L_{k/\Lambda}) = \mathfrak m_R/\mathfrak m_R^2$ and
$\Omega_{R/\Lambda} \otimes_R k = \Omega_{k/\Lambda}$
as the same is true for $P$.

\medskip\noindent
Finally, since $\Lambda \to P$ is smooth we see that
$\Lambda \to P$ is formally smooth by
Algebra, Proposition \href{https://stacks.math.columbia.edu/tag/00TN}{proposition smooth formally smooth (Tag 00TN)}.
Then $\Lambda \to P$ is formally smooth for the $\mathfrak m_P$-adic
topology by More on Algebra, Lemma \href{https://stacks.math.columbia.edu/tag/07EC}{lemma formally smooth (Tag 07EC)}.
This property is inherited by the completion $R$ by
More on Algebra, Lemma \href{https://stacks.math.columbia.edu/tag/07ED}{lemma formally smooth completion (Tag 07ED)}
and the proof is complete.
In fact, it turns out that whenever $\underline{R}|_{\mathcal{C}_\Lambda}$
is smooth, then $R$ is isomorphic to a completion of a smooth
algebra over $\Lambda$, but we won't use this.
\end{proof}
\end{document}
